% =============================================================================
% EXP 1: COMPARACIÓN DE VARIANTES
% =============================================================================
\subsection{EXP 1: Comparación de Variantes}

La Tabla~\ref{tab:exp1_tiempo} presenta el tiempo total promedio, desviación
estándar, mínimo y máximo para cada combinación instancia $\times$ variante
$\times$ población con 10 semillas. Los tiempos de las variantes CUDA
\textbf{incluyen tanto la ejecución de kernels como las transferencias de
memoria H$\to$D y D$\to$H}.

\subsubsection{Tiempo de ejecución total}

\begin{table}[H]
  \centering
  \caption{Tiempo total de ejecución (ms) — 10 semillas por celda}
  \label{tab:exp1_tiempo}
  \begin{tabular}{llrrr}
    \toprule
    \textbf{Instancia} & \textbf{Variante} &
    \textbf{pop=1024} & \textbf{pop=4096} & \textbf{pop=16384} \\
    \midrule
    \multirow{3}{*}{small}
        & sequential      & 567.2 $\pm$ 99.87    & 938.4 $\pm$ 103.23   & 3577.8 $\pm$ 431.03 \\
        & cuda\_basic     & 221.2 $\pm$ 127.38   & 230.0 $\pm$ 28.77    & 385.3 $\pm$ 25.97 \\
        & cuda\_optimized & 198.8 $\pm$ 26.82    & 253.9 $\pm$ 34.51    & 366.4 $\pm$ 20.29 \\
    \midrule
    \multirow{3}{*}{medium}
        & sequential      & 9189.3 $\pm$ 1578.43 & 21474.9 $\pm$ 3983.26 & 81633.9 $\pm$ 14380.83 \\
        & cuda\_basic     & 1205.9 $\pm$ 177.34  & 1278.4 $\pm$ 215.58   & 1747.8 $\pm$ 293.46 \\
        & cuda\_optimized & 1776.3 $\pm$ 263.55  & 2140.4 $\pm$ 361.20   & 2602.3 $\pm$ 443.35 \\
    \midrule
    \multirow{3}{*}{large}
        & sequential      & 143549.4 $\pm$ 8778.28 & 255574.2 $\pm$ 33000.94 & 1114741.4 $\pm$ 92439.55 \\
        & cuda\_basic     & 59904.8 $\pm$ 8408.96  & 69906.6 $\pm$ 6814.42   & 78480.4 $\pm$ 2267.68 \\
        & cuda\_optimized & 75026.9 $\pm$ 10862.60 & 103816.7 $\pm$ 10314.47 & 112651.4 $\pm$ 3252.47 \\
    \bottomrule
  \end{tabular}
\end{table}

\subsubsection{Porcentaje de soluciones factibles en la población}

\begin{table}[H]
  \centering
  \caption{Individuos factibles promedio en la población (\%)}
  \label{tab:exp1_factible}
  \begin{tabular}{llrrr}
    \toprule
    \textbf{Instancia} & \textbf{Variante} & \textbf{pop=1024} & \textbf{pop=4096} & \textbf{pop=16384} \\
    \midrule
    \multirow{3}{*}{small}
        & sequential      & 0.00  & 0.00  & 0.00  \\
        & cuda\_basic     & 22.52 & 23.93 & 23.13 \\
        & cuda\_optimized & 22.52 & 23.93 & 23.13 \\
    \midrule
    \multirow{3}{*}{medium}
        & sequential      & 0.00  & 0.00  & 0.00  \\
        & cuda\_basic     & 11.64 & 11.94 & 12.88 \\
        & cuda\_optimized & 11.64 & 11.94 & 12.88 \\
    \midrule
    \multirow{3}{*}{large}
        & sequential      & 0.00  & 0.00  & 0.00  \\
        & cuda\_basic     & 5.43  & 5.60  & 5.64  \\
        & cuda\_optimized & 5.43  & 5.60  & 5.64  \\
    \bottomrule
  \end{tabular}
\end{table}

\subsubsection{Calidad de solución (mejor fitness)}

\begin{table}[H]
  \centering
  \caption{Mejor fitness promedio — instancia \texttt{large}}
  \label{tab:exp1_fitness}
  \begin{tabular}{lrrr}
    \toprule
    \textbf{Variante} & \textbf{pop=1024} & \textbf{pop=4096} & \textbf{pop=16384} \\
    \midrule
    sequential        & 0.207228 & 0.207512 & 0.209281 \\
    cuda\_basic       & 0.193627 & 0.195265 & 0.196586 \\
    cuda\_optimized   & 0.193627 & 0.195265 & 0.196586 \\
    \bottomrule
  \end{tabular}
\end{table}
\noindent\textit{Nota: ambas variantes CUDA obtienen el mismo fitness promedio
porque comparten el mismo kernel de reproducción y reparación; la optimización
solo afecta la evaluación de fitness y las transferencias, no la lógica
evolutiva subyacente.}

\subsubsection{Desglose de tiempos CUDA (instancia large, pop 4096)}

Los tiempos reportados a continuación corresponden a los totales acumulados
durante las 300 generaciones. Nótese que la transferencia H$\to$D ocurre una
sola vez al inicio (subida de datos de instancia), mientras que D$\to$H ocurre
cada generación para los escalares de fitness, penalización y validez.

\begin{table}[H]
  \centering
  \caption{Desglose de tiempos CUDA (ms totales, 300 generaciones, large, pop 16384)}
  \label{tab:exp1_cuda}
  \begin{tabular}{lrrrr}
    \toprule
    \textbf{Variante} &
    \textbf{Kernel fitness} &
    \textbf{Kernel reproducción} &
    \textbf{H$\to$D} &
    \textbf{D$\to$H} \\
    \midrule
    cuda\_basic      & 510.99 & 76634.07 & 6.41 & 21.04 \\
    cuda\_optimized  & 516.24 & 110790.40 & 6.39 & 17.64 \\
    \bottomrule
  \end{tabular}
\end{table}
% =============================================================================
% EXP 2: EFECTO DEL TAMAÑO DE BLOQUE
% =============================================================================
\subsection{EXP 2: Efecto del Tamaño de Bloque}

\begin{table}[H]
  \centering
  \caption{Tiempo de ejecución (ms) según tamaño de bloque — instancia \texttt{large}, pop 16384, 10 semillas}
  \label{tab:exp2}
  \begin{tabular}{lrrrr}
    \toprule
    \textbf{Variante} & \textbf{bs=32} & \textbf{bs=64} & \textbf{bs=128} &
    \textbf{bs=256} \\
    \midrule
    cuda\_basic      & 1712.5 & 1712.7 & 1772.1 & 1909.0 \\
    cuda\_optimized  & 2578.8 & 2593.7 & 2574.5 & 2563.7 \\
    \bottomrule
  \end{tabular}
\end{table}
\textit{Nota: Los tiempos de EXP 2 se obtuvieron con una configuración
reducida (100 generaciones) para acelerar el barrido de bloques. El tamaño
óptimo para \texttt{cuda\_basic} es bs=32 o bs=64, y para
\texttt{cuda\_optimized} es bs=256. Las diferencias son marginales ($\pm$10\%).}

% =============================================================================
% EXP 3: SPEED-UP GPU vs. CPU
% =============================================================================
\subsection{EXP 3: \emph{Speed-up} GPU vs.\ CPU}

El \emph{speed-up} se calcula como $T_{\text{seq}} / T_{\text{GPU}}$, donde
$T_{\text{GPU}}$ incluye el tiempo total de la variante CUDA (kernels +
transferencias H$\to$D y D$\to$H).

\begin{table}[H]
  \centering
  \caption{\emph{Speed-up} promedio respecto de la variante secuencial}
  \label{tab:exp3}
  \begin{tabular}{llrr}
    \toprule
    \textbf{Instancia} & \textbf{Población} &
    \textbf{Speed-up CUDA Basic} & \textbf{Speed-up CUDA Opt.} \\
    \midrule
    \multirow{3}{*}{small}
        & 1024  & 2.56$\times$ & 2.85$\times$ \\
        & 4096 & 4.08$\times$ & 3.70$\times$ \\
        & 16384 & 9.29$\times$ & 9.76$\times$ \\
    \midrule
    \multirow{3}{*}{medium}
        & 1024  & 7.62$\times$ & 5.17$\times$ \\
        & 4096 & 16.80$\times$ & 10.03$\times$ \\
        & 16384 & 46.71$\times$ & 31.37$\times$ \\
    \midrule
    \multirow{3}{*}{large}
        & 1024  & 2.40$\times$ & 1.91$\times$ \\
        & 4096 & 3.66$\times$ & 2.46$\times$ \\
        & 16384 & 14.20$\times$ & 9.90$\times$ \\
    \bottomrule
  \end{tabular}
\end{table}
\noindent\textit{El mayor speed-up se alcanza en la instancia \texttt{medium}
con población 16384 ($46.71\times$ para CUDA Basic), donde la carga
computacional justifica ampliamente el uso de GPU. En \texttt{large} con
población 16384 se obtiene un speed-up de $14.20\times$.}

% =============================================================================
% EXP 4: EFECTO DEL TAMAÑO DE POBLACIÓN
% =============================================================================
\subsection{EXP 4: Efecto del Tamaño de Población}

\begin{table}[H]
  \centering
  \caption{Tiempo y fitness según tamaño de población — instancia \texttt{large}}
  \label{tab:exp4}
  \begin{tabular}{llrrr}
    \toprule
    \textbf{Variante} & \textbf{Métrica} &
    \textbf{pop=1024} & \textbf{pop=4096} & \textbf{pop=16384} \\
    \midrule
    \multirow{2}{*}{sequential}
        & Tiempo (ms) & 143549.4 & 255574.2 & 1114741.4 \\
        & Fitness     & 0.207228 & 0.207512 & 0.209281 \\
    \midrule
    \multirow{2}{*}{cuda\_basic}
        & Tiempo (ms) & 59904.8  & 69906.6  & 78480.4 \\
        & Fitness     & 0.193627 & 0.195265 & 0.196586 \\
    \midrule
    \multirow{2}{*}{cuda\_optimized}
        & Tiempo (ms) & 75026.9  & 103816.7 & 112651.4 \\
        & Fitness     & 0.193627 & 0.195265 & 0.196586 \\
    \bottomrule
  \end{tabular}
\end{table}
\noindent\textit{El fitness mejora monótonamente al aumentar la población. En
GPU el tiempo escala de forma sub-lineal: de pop 1024 a 16384 el tiempo crece
solo 1.3$\times$ en CUDA Basic (de 60\,s a 78\,s), mientras que en CPU se
multiplica por 7.8$\times$ (de 143\,s a 1114\,s).}